\item La \textit{ley de atmósferas} establece que la densidad de número de moléculas en la atmósfera depende de la altura $y$ sobre el nivel del mar de acuerdo a
$$ n_V(y) = n_0 e^{-m_0gy/k_BT} $$
donde $n_0$ es la densidad de número al nivel del mar ($y = 0$). La altura promedio de una molécula en la atmósfera de la Tierra está dada por
$$ y_{prom} = \frac{\int_0^\infty y n_V(y) dy}{\int_0^\infty n_V(y) dy} = \frac{\int_0^\infty y e^{-m_0gy/k_BT} dy}{\int_0^\infty e^{-m_0gy/k_BT} dy} $$
(a) Pruebe que esta altura promedio es igual a $k_BT/m_0g$. (b) Evalúe la altura promedio, si supone que la temperatura es $10^\circ\text{C}$ y la masa molecular es $28.9\text{ u}$, uniformes a través de la atmósfera.
